
\documentclass[12pt]{amsart}
\usepackage{geometry} % see geometry.pdf on how to lay out the page. There's lots.
\geometry{a4paper} % or letter or a5paper or ... etc
\usepackage[T1]{fontenc}
\usepackage[latin9]{inputenc}
\usepackage{amsmath}
\usepackage{amsaddr}
\usepackage{hyperref}
\usepackage{dirtytalk}
\usepackage{float}
\usepackage{listings}
\usepackage{color}
\usepackage{tikz}
\usepackage[shortlabels]{enumitem}

 
\definecolor{codegreen}{rgb}{0,0.6,0}
\definecolor{codegray}{rgb}{0.5,0.5,0.5}
\definecolor{stringcolor}{rgb}{0.7,0.23,0.36}
\definecolor{backcolour}{rgb}{0.95,0.95,0.92}
\definecolor{keycolor}{rgb}{0.007,0.01,1.0}
\definecolor{itemcolor}{rgb}{0.01,0.0,0.49}
 
\lstdefinestyle{mystyle}{
    %backgroundcolor=\color{backcolour},   
    commentstyle=\color{codegreen},
    keywordstyle=\color{keycolor},
    numberstyle=\tiny\color{codegray},
    stringstyle=\color{stringcolor},
    basicstyle=\footnotesize,
    breakatwhitespace=false,         
    breaklines=true,                 
    captionpos=b,                    
    keepspaces=true,                 
    numbers=left,                    
    numbersep=5pt,                  
    showspaces=false,                
    showstringspaces=false,
    showtabs=false,                  
    tabsize=2
}
 
\lstset{style=mystyle}

\lstdefinelanguage{Swift}{
  keywords={associatedtype, class, deinit, enum, extension, func, import, init, inout, internal, let, operator, private, protocol, public, static, struct, subscript, typealias, var, break, case, continue, default, defer, do, else, fallthrough, for, guard, if, in, repeat, return, switch, where, while, as, catch, dynamicType, false, is, nil, rethrows, super, self, Self, throw, throws, true, try, associativity, convenience, dynamic, didSet, final, get, infix, indirect, lazy, left, mutating, none, nonmutating, optional, override, postfix, precedence, prefix, Protocol, required, right, set, Type, unowned, weak, willSet},
  ndkeywords={class, export, boolean, throw, implements, import, this},
  sensitive=false,
  comment=[l]{//},
  morecomment=[s]{/*}{*/},
  morestring=[b]',
  morestring=[b]"
}

\lstset{emph={Int,count,abs,repeating,Array}, emphstyle=\color{itemcolor}}


\title{Week 09}

\date{\today}

\lstset{style=mystyle}

%%% BEGIN DOCUMENT
\begin{document}
\maketitle


Name: 
\\

Collaborators (if any): 

\section{Tasks}

\subsection{True or False}
Choose T or F for each of the following statements and briefly explain why.  
\begin{enumerate}
\item T/F Searching in a skip list takes $\Theta (\log n)$ time with high probability, but
could take $\Omega (2n)$ time with nonzero probability.  (Note: this statement could actually be answered as either true or false with proper justification) 
\item To implement a dictionary, you should always use a hash table over an AVL tree if your main priority is speed.
\item The number of collisions in a hash table is solely dependent on the table capacity and the hash function.
\end{enumerate}

\section{Exercises/Problems}

\subsection{Calculate Hash Values for Given Keys} 
Given the following input (4322, 1334, 1471, 9679, 1989, 6171, 6173, 4199) and the hash function x mod 10, which of the following statements are true? 
\begin{enumerate}[i.]
\item 9679, 1989, 4199 hash to the same value
\item 1471, 6171 hash to the same value
\item All elements hash to the same value
\item Each element hashes to a different value 
\end{enumerate}
\begin{enumerate}[(A)]
\item i only
\item ii only 
\item i and ii only
\item iii or iv 
\end{enumerate}

\subsection{Given a hash table with keys, verify/find possible sequence of keys}
For a given hash table, we can verify which sequence of keys can lead to that hash table. However, to find possible sequences leading to a given hash table, we need to consider all possibilities. A hash table of length 10 uses open addressing with hash function h(k)=k mod 10, and linear probing. After inserting 6 values into an empty hash table, the table is as shown below. 

\begin{tabular}{| l | r |}
\hline
  0 &   \\
  1 &   \\
  2 & 42  \\
  3 & 23 \\
  4 & 34 \\
  5 & 52  \\
  6 & 46  \\
  7 & 33 \\
  8 &  \\
  9 &  \\

\hline
\end{tabular}
\\

Which one of the following choices gives a possible order in which the key values could have been inserted in the table? 
\begin{enumerate}[(A)]
\item 46, 42, 34, 52, 23, 33
\item 34, 42, 23, 52, 33, 46 
\item 46, 34, 42, 23, 52, 33
\item 42, 46, 33, 23, 34, 52
\end{enumerate}


\end{document}
